\begin{answer}
    \begin{enumerate}
        \item The learning rate will not ultimately effect whether we would converge, as it would either lead to bigger or smaller gradient updates but not lead to a finite minimum in the loss landscape. 
    
        \item Same reasoning as above, the learning rate doesn't affect whether we have a finite minimum to converge to. It might help us practically to "converge" because the parameter updates will be small enough so that our learning algorithm will stop (if difference of parameters \(< \epsilon\)).
    
        \item If we scale our parameters by e.g. 1/n, our training algorithm will just need b times as long to reach a specific magnitude of parameters, and this will not influence the loss landscape and having a convergable minimum.
    
        \item A regularization term \(\|\|\theta\|\|^2\) would most definitely help our learning algorithm to converge, as the penalty induced would grow the bigger, the more advanced our convergence process is. There will be a certain point, when the regularization term and fit loss will balance out when we reached a finite minimum.
    
        \item This might cause the data points/labels to spread out in a way that makes it linearly separable and our training algorithm to converge to a finite minimum, but it might cause the opposite based on noise being inherently random.
    \end{enumerate}
    
    \vspace{2em}
\end{answer}
